% Part: computability
% Chapter: recursive-functions
% Section: normal-form

\documentclass[../../../include/open-logic-section]{subfiles}

\begin{document}

\olfileid{cmp}{rec}{nft}
\olsection{Teorema Wangun Normal}

\begin{thm}[Teorema Wangun Normal Kleene]
\ollabel{thm:kleene-nf}
Ana relasi rekursif primitif $T(e, x, s)$ lan
fungsi rekursif primitif $U(s)$, kanthi sipat iki:
yen $f$ fungsi rekursif parsial sembarang,
mula kanggo sawenehing~$e$,
\[
f(x) \simeq U(\umin{s}{T(e, x, s)})
\]
kanggo saben $x$.
\end{thm}

\begin{explain}
Bukti teorema wangun normal iki rumit,
nanging gagasan dhasare prasaja. Saben
fungsi rekursif parsial nduwé \emph{indeks}~$e$,
yaiku, sacara intuitif, wilangan kang ngode
program utawa definisine. Yen $f(x)
\fdefined$, komputasi iku bisa dicathet
kanthi sistematis lan dikode nganggo
sawenehing wilangan~$s$; kasunyatan yen $s$
ngode komputasi~$f$ kanggo input~$x$
bisa dipriksa kanthi cara rekursif primitif
mung nganggo $x$ lan definisi~$e$.
Mula relasi~$T$, kang nyatakake 'fungsi
mawa indeks~$e$ nduwé komputasi kanggo
input~$x$, lan $s$ ngode komputasi iki,'
iku rekursif primitif. Yen cathetan
komputasi~$s$ wis pepak, asile
saka~$s$ yaiku nilai~$f(x)$, lan
nilai iki uga bisa diolehake saka~$s$
kanthi cara rekursif primitif.

Teorema wangun normal nuduhake yen
kanggo netepake fungsi rekursif parsial
sembarang cukup mung siji panggolekan
tanpa wates. Sejatine kita bisa nlusuri
kabeh wilangan nganti nemokake wilangan
kang ngode komputasi fungsi mawa indeks~$e$
kanggo input~$x$. Kita bisa nggunakake
wilangan~$e$ minangka 'jeneng' fungsi
rekursif parsial, lan nulis $\cfind{e}$
kanggo fungsi~$f$ kang ditegesi dening
persamaan ing teorema mau. Gatekna yen
fungsi rekursif parsial sembarang bisa
nduwé luwih saka siji indeks---malah
saben fungsi rekursif parsial nduwé
indeks tanpa wates cacahé.
\end{explain}
\end{document}
